% index.tex -- index of Module 6 multiple-choice questions in the 2019-2025 NSW HSC Physics papers.
% Descriptions are paraphrases, not question text.
\documentclass[12pt]{article}
\usepackage{quiz}
\usepackage{longtable}
\hypersetup{pdftitle={HSC Physics Module 6 Multiple-Choice Questions 2019-2025 - Index}}
\begin{document}
\quiztitle{Year 12 Physics}{HSC Module 6 Multiple-Choice Questions 2019–2025: Index}
Every Section I question from the 2019--2025 NSW HSC Physics papers that tests Module 6, with the
answer from the NESA marking guidelines. Topics: T1 charged particles in electric and magnetic fields;
T2 the motor effect; T3 electromagnetic induction; T4 applications of the motor effect. M7/M8 = also uses
Module 7 or 8 ideas. The descriptions are short summaries, not the question text.

{\small\renewcommand\arraystretch{1.3}
\begin{longtable}{|c|c|c|p{108mm}|c|}
  \hline
  \textbf{Year} & \textbf{Qu} & \textbf{Topic} & \textbf{What it tests} & \textbf{Ans}\\ \hline
  \endhead
    2019 & 5 & T3 & Two coils on an iron rod: galvanometer when the primary switch opens & D\\ \hline
    2019 & 7 & T3 & Lenz's law: induced current and coil polarity as a magnet moves away & D\\ \hline
    2019 & 16 & T1 & Charged particle in an E field plus gravity: field for equal upward deflection & D\\ \hline
    2019 & 17 & T1 & Cathode ray beside a current-carrying wire: spot shift when the current drops & D\\ \hline
    2019 & 18 & T3 & Battery-lamp loop moving through a B-field region: brightness entering and leaving & B\\ \hline
    2020 & 2 & T4 & AC induction motor: the part it does not need (brushes) & A\\ \hline
    2020 & 7 & T4 & Matching a rectified voltage graph to a split-ring DC generator & D\\ \hline
    2020 & 10 & T1 & Crossed E and B fields: directions for an undeflected electron & B\\ \hline
    2020 & 14 & T2 & Parallel wires: force after one current is doubled and reversed & C\\ \hline
    2020 & 19 & T3 & Rod spinning in a plane containing B: induced emf is zero & C\\ \hline
    2021 & 2 & T1 & Direction of the electric force on a moving positive charge & C\\ \hline
    2021 & 7 & T3 & Ideal transformer: secondary current 4 times primary gives type and turns ratio & C\\ \hline
    2021 & 10 & T4 & Eddy currents: force on a copper block as a magnet sweeps past & B\\ \hline
    2021 & 12 & T4 & DC motor: back emf against angular velocity (linear) & D\\ \hline
    2021 & 17 & T2 & Hanging parallel conductors carrying AC: comparing spring-balance readings & C\\ \hline
    2021 & 18 & T1 & Charged mass falling between plates: path after the field is switched off & A\\ \hline
    2021 & 20 & T1/\allowbreak M8 & Photoelectron circling in a B field: radius when photon energy doubles & D\\ \hline
    2022 & 1 & T3 & Ideal transformer: number of secondary turns from voltages & C\\ \hline
    2022 & 3 & M8/\allowbreak T1 & Identifying radiation type from its deflection between charged plates & D\\ \hline
    2022 & 4 & T2 & Motor effect: force direction on a wire between magnet poles & C\\ \hline
    2022 & 12 & T1 & Velocity selector: velocity direction of an undeflected positive charge & D\\ \hline
    2022 & 15 & T2 & Parallel wires: force when separation is 4 times and both currents doubled & A\\ \hline
    2022 & 18 & T1/\allowbreak M8 & Millikan oil drop: force balance with the field off and on & D\\ \hline
    2022 & 19 & T4 & Hand generator with two bulbs in series: easier to turn with less current & A\\ \hline
    2023 & 2 & T3 & Power transmission: step-up, high voltage and low current, then step-down & C\\ \hline
    2023 & 10 & T4 & Rotated current loop: force on one side unchanged, torque decreases & D\\ \hline
    2023 & 12 & T1 & Plates moved apart at the same voltage: kinetic energy gained over a fixed distance & B\\ \hline
    2023 & 16 & T2 & Rod levitated above another rod: current directions and support force & C\\ \hline
    2023 & 18 & T1 & Mass increased: projectile path in gravity vs a charge's path in an E field & A\\ \hline
    2023 & 19 & T2/\allowbreak M7 & Relativity explanation of the force between wires: apparent charge density & B\\ \hline
    2024 & 4 & T4 & DC motor needs a DC supply and split-ring commutator to keep rotating & D\\ \hline
    2024 & 8 & T3 & Ideal transformer: change the primary turns to double the output voltage & B\\ \hline
    2024 & 10 & T2 & F = lIB sin($\theta$): rotation angle that halves the force on a rod & C\\ \hline
    2024 & 15 & T3 & Rod spinning in a plane perpendicular to B: constant emf & D\\ \hline
    2024 & 17 & T1 & Cyclotron: where the particle accelerates & C\\ \hline
    2024 & 18 & T3 & Lenz's law and mutual induction: arrangement giving the same current direction & D\\ \hline
    2024 & 19 & T1 & Velocity selector entered at 2v: change in size and direction of acceleration & C\\ \hline
    2025 & 2 & T3 & Step-up transformer: type and turns from 240 V to 2200 V & A\\ \hline
    2025 & 6 & T1 & Electron curving in a B field: identify the field direction & A\\ \hline
    2025 & 11 & T4 & DC motor slowed by a load: back emf falls, so current rises & D\\ \hline
    2025 & 14 & T1 & Proton entering a field: field type that raises kinetic energy (electric) & B\\ \hline
    2025 & 17 & T3 & Loop moving away from a straight wire: emf shape and direction & C\\ \hline
\end{longtable}}

Summary: 42 questions, 5 to 7 in each paper (T1: 15, T3: 12, T4: 8, T2: 7). Transformers appear almost
every year; velocity selectors, charges in an electric field with gravity, Lenz's law, the force between
parallel wires and back emf come up often.
\end{document}
